\chapter{El \emph{string interner}}
\label{cap:interner}

\begin{intuicion}
Las identidades del grafo (\texttt{"alice"}, \texttt{"db-01"},
\texttt{"10.0.0.42"}) son \emph{strings} para los humanos, pero
manipularlas como \texttt{String} en el matcher es caro: cada
\texttt{Relation} arrastra dos punteros al heap, cada comparación
recorre los caracteres byte por byte, cada hash recalcula desde
cero. El \emph{interner} resuelve esto canónicamente: cada string
única se almacena \emph{una vez} y se referencia por un
\texttt{StrId} de 4~bytes (\texttt{u32}). El matcher pasa
\texttt{StrId}s en su \emph{hot loop}; sólo al \emph{emit} de un
resultado se resuelve el \texttt{StrId} de vuelta a la string
visible al humano.
\end{intuicion}

\section{El problema}

Consideremos un grafo de 10\,M aristas donde cada arista guarda dos
strings (\texttt{source\_id}, \texttt{dest\_id}) de 16 caracteres en
promedio.

\begin{itemize}
    \item Como \texttt{String}: cada string es \(16\)~B en heap +
    \(24\)~B de \emph{header} (puntero, longitud, capacidad). Total
    por string: \(40\)~B. Por arista: \(80\)~B sólo en strings.
    \item Total: \(10^7 \times 80 = 800\)~MB sólo para los IDs de
    las aristas.
\end{itemize}

\noindent
Pero las identidades se repiten masivamente: en un grafo típico hay
\(\sim 100\)\,K usuarios distintos, \(\sim 10\)\,K hosts distintos,
\(\sim 1\)\,M IPs distintas. La cardinalidad real del conjunto de
strings es del orden de \(10^6\), no \(10^7\); guardar la string
completa en cada arista es \(\sim 10\times\) más memoria de lo
necesario.

Además, el matcher hace \(\sim 10^9\) comparaciones de IDs por
hunt (cada paso del DFS chequea ``¿este vértice ya está bound?''
contra los bindings actuales). Comparar dos \texttt{String}s de 16~B
es 16 cargas de memoria + 16 comparaciones de byte. Comparar dos
\texttt{u32} es una sola instrucción.

\section{Definición}

\begin{definition}[String interner]
Un \emph{interner} sobre un alfabeto \(\Sigma\) es una estructura de
datos que mantiene:
\begin{itemize}
    \item Un mapeo inyectivo \(\sigma \colon \Sigma \to \mathbb{N}\)
    (la \emph{intern function}).
    \item Una resolución \(\sigma^{-1} \colon \mathbb{N} \to \Sigma\)
    (la \emph{resolve function}).
\end{itemize}
Las dos operaciones soportadas son:
\begin{itemize}
    \item \texttt{intern}\((s)\): si \(s\) ya está internada, devuelve
    su \(\mathrm{id}\); si no, le asigna el próximo \(\mathrm{id}\)
    libre y la guarda.
    \item \texttt{resolve}\((\mathrm{id})\): devuelve la string
    asociada al \(\mathrm{id}\).
\end{itemize}
\end{definition}

\section{Implementación: tabla bidireccional}

La realización canónica:

\begin{lstlisting}[language=Rust]
pub struct StringInterner {
    forward: HashMap<Arc<str>, StrId>,  // string  -> id
    backward: Vec<Arc<str>>,            // id      -> string
}

#[derive(Copy, Clone, PartialEq, Eq, Hash)]
pub struct StrId(pub u32);

impl StringInterner {
    pub fn intern(&mut self, s: &str) -> StrId {
        if let Some(id) = self.forward.get(s) {
            return *id;
        }
        let arc: Arc<str> = Arc::from(s);
        let id = StrId(self.backward.len() as u32);
        self.backward.push(arc.clone());
        self.forward.insert(arc, id);
        id
    }

    pub fn resolve(&self, id: StrId) -> Arc<str> {
        self.backward[id.0 as usize].clone()
    }
}
\end{lstlisting}

\noindent
Tres detalles importantes:

\begin{enumerate}
    \item \textbf{\texttt{Arc<str>} en ambos lados.} Tanto la
    \texttt{HashMap} como el \texttt{Vec} guardan
    \texttt{Arc<str>} ---no \texttt{String}---. Esto permite que
    \texttt{resolve} devuelva un \texttt{Arc<str>} clonado (refcount
    bump, $\sim$5~ns) en vez de una \texttt{String} alocada de
    nuevo. Es la pieza que destrabó el \(55\%\) de mejora en los
    benches cap-saturados (capítulo~\ref{cap:hunt-result}).
    \item \textbf{\texttt{StrId} es \texttt{u32}, no \texttt{u64}.}
    Asume \(< 2^{32} = 4.3 \cdot 10^9\) strings distintas. En
    GraphHunter eso es muchísimo: el grafo más grande visto en
    producción tiene \(\sim 10^7\) strings únicas, tres órdenes de
    magnitud por debajo del techo.
    \item \textbf{No requiere \texttt{remove}.} El interner es
    \emph{append-only}; los IDs son estables por la vida de la
    sesión. Cuando la sesión se cierra, el interner se descarta
    completo. Esa propiedad simplifica todo: no hay
    \emph{tombstones}, no hay re-asignación, no hay locks complejos
    para ``hueco liberado''.
\end{enumerate}

\section{Análisis de costos}

\subsection{\texttt{intern}: amortizado \(O(1)\)}

\begin{itemize}
    \item Lookup en la \texttt{HashMap}: \(O(1)\) esperado (hash de
    \(s\), comparación contra el bucket). Para \(s\) de longitud
    \(L\), el hash cuesta \(O(L)\) ---no es independiente del
    tamaño de la string---. Pero \(L\) es chico ($\sim$10--40
    bytes para identidades típicas) y constante en escala
    asintótica.
    \item Si miss: alocación de \texttt{Arc<str>}, \texttt{push} al
    \texttt{Vec}, \texttt{insert} al \texttt{HashMap}. Los tres son
    \(O(1)\) amortizado (la \texttt{Vec} y la \texttt{HashMap}
    crecen geométricamente).
\end{itemize}

\subsection{\texttt{resolve}: \(O(1)\)}

Indexación directa al \texttt{Vec}. \texttt{Arc::clone} es un
\texttt{fetch\_add} atómico ($\sim$5 ns). Sin alocación, sin
copia de bytes.

\subsection{Espacio}

Sea \(N\) el número de strings únicas, \(L\) la longitud promedio.
La memoria total es:

\[
\underbrace{N \cdot L}_{\text{bytes de las strings}} +
\underbrace{N \cdot 16}_{\text{Arc<str> headers en backward}} +
\underbrace{N \cdot (\sim 40)}_{\text{HashMap buckets}}
\approx N \cdot (L + 56).
\]

Para \(N = 10^6\), \(L = 16\): \(\sim 72\)~MB. Comparado con los
\(800\)~MB que costaría almacenar las strings repetidas en cada
arista, el ahorro es \(\sim 11\times\).

\subsection{Comparaciones rápidas}

Una vez internadas dos strings, \(\sigma(s_1) = \sigma(s_2) \iff
s_1 = s_2\). Comparar \texttt{StrId == StrId} es una sola
comparación de \texttt{u32}; \texttt{HashMap<StrId, V>} usa el
\texttt{u32} mismo como hash (no hay re-hash de bytes). Esto es lo
que hace que el \emph{hot loop} del DFS pueda chequear
``¿\texttt{StrId} \(v\) ya está bound?'' en una sola comparación
de palabra de máquina.

\section{Variantes de diseño}

\subsection{Por qué \emph{no} usamos \texttt{Symbol} de string-interner crate}

Hay crates de Rust (\texttt{string-interner},
\texttt{lasso}) que hacen lo mismo. Las consideramos. Los rechazamos
por dos razones:

\begin{itemize}
    \item Devuelven \texttt{\&str} con \emph{lifetime} atado al
    interner. El motor necesita \texttt{Arc<str>} (ownership compartido,
    serialización por refcount). Un \texttt{\&str} se vuelve estorbo
    cuando el resultado del hunt cruza un \emph{thread boundary} o
    se serializa al frontend.
    \item Soportan \emph{dropping} de strings (con
    \emph{tombstones}, \emph{garbage collection}). En GraphHunter
    nunca se borra; pagar el overhead de soportar drops sería
    desperdicio.
\end{itemize}

\subsection{Por qué \emph{no} usamos un \emph{trie} de prefix sharing}

Un \emph{trie} podría compartir prefijos comunes
(\texttt{"CORP\textbackslash{}m.gomez"} y
\texttt{"CORP\textbackslash{}j.lopez"} compartirían 5 bytes). El
ahorro teórico es real pero dos factores juegan en contra:

\begin{itemize}
    \item Lookup en trie es \(O(L)\) con \(L\) cargas de memoria
    no-secuenciales. Para \(L = 16\), eso son 16 cache misses
    posibles ---\(>200\) ciclos. Hash table (con buen hash) es \(\sim
    10\) ciclos.
    \item Devolver \texttt{Arc<str>} desde un trie requiere
    materializar la string concatenando los prefijos. La
    concatenación cuesta lo que ahorraste.
\end{itemize}

\noindent
La hash table \texttt{Arc<str>} gana en práctica.

\section{Concurrencia}

El \texttt{StringInterner} de GraphHunter es \texttt{\&mut self}
para \texttt{intern}: los writers (el ingestor) toman acceso
exclusivo durante la inserción. Los \texttt{StrId} resueltos son
estables, así que los readers (el matcher) pueden trabajar con
copias de los IDs sin tocar el interner.

\begin{enelcodigo}
\filepath{core/graph-engine/src/interner.rs}: \texttt{StringInterner}
con \texttt{intern}, \texttt{resolve}. Bajo concurrencia
ingest+matching, el interner se envuelve en un \texttt{RwLock}; el
ingestor toma write-lock para inserts, los matchers toman
read-lock para resolves. Los matchers nunca bloquean a otros
matchers ---el bottleneck efectivo es el ingest, que típicamente es
sigle-threaded.
\end{enelcodigo}

\section{Por qué importa para el resto del libro}

El \texttt{StrId} es la representación interna de una identidad en
todo el motor:

\begin{itemize}
    \item \texttt{CompactRelation} (capítulo~\ref{cap:compact-relation})
    almacena \texttt{source\_sid: StrId} y \texttt{dest\_sid: StrId},
    cada uno 4~B en lugar de los 24~B de un \texttt{String}.
    \item NLF (capítulo~\ref{cap:nlf-detalle}) indexa por
    \texttt{StrId} para tener \(O(1)\) lookup.
    \item k-hop reachability (capítulo~\ref{cap:khop-detalle}) usa
    \texttt{StrId} como índice en sus arrays \texttt{bits} y
    \texttt{rank}.
    \item El DFS \emph{hot loop} pasa \texttt{StrId} entre frames de
    la recursión sin tocar strings.
    \item Sólo el \emph{emit} de un resultado al fin del DFS resuelve
    \texttt{StrId} $\rightarrow$ \texttt{Arc<str>}.
\end{itemize}

Esto último importa: la \emph{cantidad} de resoluciones es
proporcional al número de resultados (\(\leq 10\,000\) por hunt en
default), no al trabajo del DFS (\(10^9\) comparaciones). Pagar
\(5\)~ns por resolución es 50~$\mu$s de overhead total ---trivial.

\section{Resumen}

\begin{itemize}
    \item Interner: \texttt{intern}\((s) \to \mathrm{id}\) y
    \texttt{resolve}\((\mathrm{id}) \to s\). Tabla bidireccional
    (\texttt{HashMap} + \texttt{Vec}).
    \item \texttt{StrId} es \texttt{u32} (4~B); \texttt{Arc<str>}
    en el backing store permite \texttt{resolve} con refcount bump,
    sin alocar.
    \item Memoria: \(\sim 11\times\) más eficiente que repetir las
    strings en cada \texttt{Relation}; comparaciones \(\sim 16\times\)
    más rápidas que comparar strings byte a byte.
    \item Ningún \texttt{drop}; \texttt{StrId}s estables por la vida
    de la sesión.
\end{itemize}

\begin{ejercicios}
\begin{enumerate}
    \item \starred{} Calculá la memoria que ahorraría el interner en
    un grafo de \(10^8\) aristas con \(N = 10^6\) strings únicas de
    longitud promedio \(L = 20\)~B. Compará con almacenar la string
    completa en cada arista.

    \item ¿Por qué la \texttt{intern} hash function NO debería ser
    crypto-grade (SHA-256)? ¿Qué propiedades sí necesita? (Pista:
    pensá adversarial: ¿qué pasa si un atacante puede elegir las
    strings que entran al interner?)

    \item \starred{} Implementá una variante del interner que NO
    use \texttt{HashMap} para \texttt{forward}, sino un \emph{Vec}
    ordenado con búsqueda binaria. ¿Qué cambia en la complejidad de
    \texttt{intern}? ¿Cuándo conviene?

    \item Imaginá que mañana la red corporativa explota y el
    interner crece a \(N = 10^9\) strings. \texttt{StrId = u32} ya
    no alcanza. ¿Qué cambia? Discutí: ¿migrar a \texttt{u64}?
    ¿segmentar por \texttt{EntityType}? ¿ambas?

    \item ¿Por qué el interner usa \texttt{Arc<str>} en
    \texttt{backward} y no \texttt{String}? ¿Qué cambia para
    \texttt{resolve}?

    \item En un sistema con concurrencia ingest + matching, el
    matcher tiene un \texttt{StrId} y quiere resolverlo a la string,
    pero el ingestor está ejecutando \texttt{intern} en paralelo.
    ¿Hay riesgo de \emph{torn read}? ¿Qué garantiza la consistencia?
\end{enumerate}
\end{ejercicios}
